\documentclass[letterpaper,11pt]{moderncv}
\moderncvstyle{banking}
\moderncvcolor{black}

\usepackage[letterpaper, margin=0.8in]{geometry}
\usepackage{xcolor}
\usepackage{fontspec}
\setmainfont{Times New Roman}
\usepackage[unicode]{hyperref}

\name{Saeyoung}{Kim}
\address{Cambridge, MA, USA}{}{}
\phone[mobile]{+1 (351) 322 5510}
\email{saeyoung\_kim@uml.edu}
\social[linkedin]{saeyoung-macx-kim}

\recipient{Hiring Team}{Busek Co. Inc.\\Natick, MA}
\date{\today}
\opening{Dear Hiring Team,}
\closing{Sincerely,}

\begin{document}

\makelettertitle

I'm applying for the Senior Electrical Test Engineer position. Testing and validating electronic hardware for spaceflight is what I've been doing, and Busek's work on electric propulsion is the kind of program I want to keep doing it on.

At Hanwha Systems, I led test campaigns for military SAR satellite programs across EM, QM, and FM builds. I designed the EGSE architecture, configured test setups, and spent most of my time in the lab with oscilloscopes, DMMs, and power supplies, running system-level tests and tracking down anomalies. When hardware failed during integration, I did the root cause analysis, worked through schematics and firmware logs to isolate the issue, and coordinated fixes with the design teams. I also wrote the controlled test procedures and documentation used for both internal engineering reviews and customer deliverables, and supported environmental testing campaigns including thermal vacuum and vibration qualification. Part of the role was mentoring junior engineers on test methodology and lab practices.

My PhD was seven years of building and testing electronics from scratch. I designed circuits, laid out PCBs, wrote firmware in C/C++ for data acquisition over UART, SPI, and I2C, and then built the test setups to characterize and qualify everything. I ran environmental and reliability tests across multiple prototype cycles, learning to find failure modes early and trace them back to root cause.

On the software side, I've been writing Python tools for automated process monitoring and anomaly detection at UMass Lowell, which has strengthened my ability to build data-driven test workflows.

Beyond the technical side, I still have strong working relationships at Hanwha Systems and SatReC (Satrec Initiative), both of which are building out satellite constellations. If Busek is looking to expand its propulsion customer base in the Korean satellite market, I'd be glad to help make those introductions.

On the ITAR front: my green card is currently pending, and I expect to have it in hand within the next few months. I'm happy to discuss timing directly. I'm based in Cambridge, so Natick is a short commute.

\makeletterclosing

\end{document}
